\chapter{讨论、结论与AI辅助说明}

\section{开放性思考与讨论（20分）}

\subsection{跨时间分布漂移导致泛化压力增大}

本实验采用历史训练、未来测试设置，真实暴露了时间漂移问题。即使训练阶段表现稳定，未来时段仍可能出现整体偏移，体现为$R^2$下降与长尾误差增加。该风险提示模型部署后可能发生性能衰减，后续可采用滚动窗口重训练、时间衰减加权与分季节建模进行缓解。

\subsection{多源异步采样带来的信息不对齐}

多源监测系统常存在采样周期差异、网络延迟与设备离线。当前采用锚表左连接与缺失填补保证了可训练性，但填补不能替代真实观测。后续可引入时间窗口匹配、缺失指示变量和按可用率分层建模，将“缺失模式”转化为可学习信号。

\subsection{异常值中信号与噪声的边界不清}

极端值可能来源于设备异常，也可能代表真实环境突变。静态物理阈值能去除明显非法值，但对边界样本仍缺乏情境化判断。下一步可结合站点历史分位数、相邻时段一致性和多源协同异常检测，提升清洗可信度。

\subsection{线性模型的可解释性与表达上限}

线性体系便于解释与教学复核，但在复杂非线性耦合场景下存在拟合上限。该限制并非纯粹缺点，它有助于区分“数据治理可改进问题”和“模型能力上限问题”，为后续迭代提供清晰方向。

\subsection{多指标联合解读的重要性}

在跨时段评估中，单看$R^2$容易产生偏差。MAE反映平均误差，RMSE对大误差更敏感，$R^2$需结合分布漂移背景解读。本文坚持多指标联合解释，避免了“只报分数、不解释机制”的常见问题。

\subsection{可复现性是实验质量的关键}

本实验采用统一脚本、统一输出目录和固定证据文件（CSV/JSON/PNG），保证结论可追溯、可复核。这一工程化链路既提升了报告可信度，也为后续迭代提供了稳定基线。

\section{结论（5分）}

本文基于真实海洋环境多源监测数据，完成了“盐度预测+环境状态解释”的完整实验流程。通过锚表融合、防泄漏控制、时间序列交叉验证与多指标联合评估，报告形成了从数据治理到模型评估、从数值预测到规律解释的闭环。

结果表明：Ridge在小时级与天级均优于普通线性回归。小时级最佳模型达到\texttt{MAE=3.6692, RMSE=8.7930}，天级最佳模型达到\texttt{MAE=6.0130, RMSE=12.7383}；聚类最优为\texttt{k=2}，轮廓系数\texttt{0.2772}，揭示了可用于误差解释的状态异质性。

总体上，本实验在课程已学方法框架内实现了完整、可复核、可扩展的分析流程，满足课程实验报告提交要求。

\section{AI辅助分析与算法改进说明（附加分，最高10分）}

本实验将AI用于“分析辅助”而非“结果替代”，主要体现在三个环节：
\begin{itemize}
    \item 协助梳理实验章节结构，使报告从代码记录升级为问题驱动叙事；
    \item 协助比较和固化线性体系内改进策略，如锚表左连接、周期特征、Ridge时序调参与防目标泄漏；
    \item 协助结果解释，将指标变化与数据机制关联，提升结论可读性。
\end{itemize}

需要说明的是，数据读取、清洗、训练、评估与可视化输出均由本地脚本完成；关键实验决策由本人根据课程要求与数据实际情况做出。AI在本实验中的角色是提升效率和表达质量的工具，这一使用方式符合课程规范中“可使用AI辅助但需说明参与环节与价值”的要求\cite{rubric_2025}。
